% 两个三角形{1,2,3}、{4,5,6}由边{3,4}相连：割点3,4，桥{3,4}；链1-2-3与三角形对照。
\begin{tikzpicture}[>=Stealth,
  v/.style={circle,draw=FigureStroke,fill=FigureBlueFill,line width=.8pt,minimum size=6.5mm,inner sep=0pt},
  vc/.style={v,draw=FigureRed,fill=FigureRedFill},
  e/.style={draw=FigureStroke,line width=.8pt},
  br/.style={draw=FigureRed,line width=1.6pt},
  note/.style={font=\FigureNoteFont\fontsize{7.5}{9}\selectfont,text=FigureMuted,align=center}]
\begin{scope}
  \node[v] (1) at (-2.1,.7) {$1$}; \node[v] (2) at (-2.1,-.7) {$2$}; \node[vc] (3) at (-.9,0) {$3$};
  \node[vc] (4) at (.5,0) {$4$}; \node[v] (5) at (1.7,.7) {$5$}; \node[v] (6) at (1.7,-.7) {$6$};
  \draw[e] (1)--(2)--(3)--(1) (4)--(5)--(6)--(4);
  \draw[br] (3)--(4);
  \node[align=center] at (-.2,-1.5) {(a) 割点$3,4$；桥$\{3,4\}$};
\end{scope}
\begin{scope}[xshift=3.3cm]
  \node[v] (1) at (0,0) {$1$}; \node[vc] (2) at (1.1,0) {$2$}; \node[v] (3) at (2.2,0) {$3$};
  \draw[br] (1)--(2)--(3);
  \node[align=center] at (1.1,-1.5) {(b) 链：割点$2$，\\两条边都是桥};
\end{scope}
\begin{scope}[xshift=6.8cm]
  \node[v] (1) at (0,.7) {$1$}; \node[v] (2) at (-.75,-.5) {$2$}; \node[v] (3) at (.75,-.5) {$3$};
  \draw[e] (1)--(2)--(3)--(1);
  \node[align=center] at (0,-1.5) {(c) 三角形：每条边\\都有替代路径};
\end{scope}
\end{tikzpicture}
